\section{Introduction}
\label{sec:intro}

A reader who wants to know which spiking graph neural network to use has, at present, only one tool:
a table of published accuracies. On Cora, that table tells them that SpikingGCN reaches $80.7$,
SpikeGraphormer $84.8$ and SpikeGCL $87.4$. The ordering is clear, the margins are small, and every
entry was produced by a different codebase, on a split the paper describes in one clause, with a
number of seeds it may or may not state. The reader has no way to know whether those three numbers
were measured on the same task.

This paper is a reproducibility and stress-testing audit of those published evaluations. Its
contribution is an executable protocol and the measurements it yields, not a new dataset or a new
model. The question is narrow: do published numbers reproduce from the released code? We took every
spiking-GNN method in our related work that declares a code repository --- three of seven; the other
four have no code to run --- executed each one's own training script in an environment pinned to its
own requirements, and changed one variable at a time.

\paragraph{What we found.} The three disagreements differ in kind, and that difference determines
what would fix each of them.

\emph{One reproduces on some datasets and not others.} SpikingGCN's released code reaches $50.0$ on
Cora over five seeds against the $80.7$ its paper reports, $64.3$ on PubMed against $77.6$, and
$70.6$ on CiteSeer against $72.5$. The CiteSeer cell reproduces to within two points, so the Cora
gap cannot be explained by a broken environment or a misconfigured loader. We re-measured five input
variants on Cora through the audited pipeline and none closes it; the ones a reader is likely to
guess make the number \emph{worse} by collapsing it to chance. We report the gap and
stop there, rather than search over input scalings until one of them produced $80.7$.

\emph{One reproduces exactly, on a split nobody else uses.} SpikeGraphormer's script, run with the
arguments in its own \texttt{run.sh}, reaches $84.7$ over three runs against a published
$84.8 \pm 0.7$: a reproduction by any standard. But the configuration uses a random class-balanced
split, and the paper's footnote defers the split to the code. Switching to the public Planetoid split
--- one flag --- gives $74.0$. Ten points of the published number are a property of the split, not of
the method, and a reader comparing it to a public-split baseline is comparing two different tasks.

\emph{One reproduces, and exposes a labelling failure underneath the field.} SpikeGCL reaches $86.6$
on Cora against $87.4 \pm 0.6$ and $77.3$ on Citeseer against $77.6$: it reproduces. But its table
cites SpikingGCN at $87.7$, and that number should not be taken at face value, because SpikingGCN's
own paper prints two values for it --- $80.7$ on its Split~I and $88.7$ on Split~II. The two later
papers we audit quote $81.5$ and $87.7$ for baseline rows from the same comparison source without saying which split they mean,
and the splits differ by five to twelve points. This is not a citation error; both numbers are
correct. It is an unlabelled choice, which is why it has survived.

\paragraph{Three independent checks on the same conclusion.} A single failed reproduction admits an
implementation explanation. Here one pipeline reproduces on CiteSeer while missing on Cora, one
reproduces exactly under a different split, and one reproduces while its own baseline table mixes
splits. No single environment, variance or citation mistake explains all three, and each can be
checked against the released scripts.

The audit also has a constructive half. Once all three methods are measured on the public split,
four accuracies become commensurable on Cora: $50.0$ (SpikingGCN), $74.0$ (SpikeGraphormer), $80.5$
(a conversion method from our companion work \citep{companion}) and $86.6$ (SpikeGCL). That ordering
differs from the published one, and it is the ordering a reader can use.

\paragraph{Contributions.}
\begin{enumerate}\itemsep2pt
\item \textbf{A pinned-environment audit of every auditable spiking-GNN method.} Three pipelines, each
run with its authors' own script and configuration; every reported value comes from the run's own
log.
\item \textbf{A taxonomy of the ways a spiking-GNN number can fail to transfer.} Dataset-dependent
non-reproduction (SpikingGCN), split-dependence within one pipeline (SpikeGraphormer), and
split-mixing across the literature (baseline rows cited at $81.5$ and $87.7$ from the two splits of
one comparison source, with no label). The
three have different remedies, and conflating them is why the problem has been invisible.
\item \textbf{A measurement of the split effect.} On Cora, changing only the split moves
SpikeGraphormer by $10.7$ points --- larger than the margin between most methods the field compares.
\item \textbf{A census, not a sample.} Three of the seven methods in our related work declare a
repository; we audited all three. The remaining four are unauditable by anyone, which is a fact about
the field rather than a limitation of this study.
\end{enumerate}

The scope is the auditable population: methods whose released code can be executed in a pinned
environment. A number can be correct for the experiment it measures and unusable for the comparison a
reader wants to make with it; the SpikeGraphormer case is exactly that, and it is the most instructive
of the three.
